% page 586 of the facsimile (image p-585 of the scan)
% block 157.5 x 254.9 mm ; left margin 8.0 mm
\pgc{-12.700mm}{0.000mm}{
\l{}
\l{}
\l{\cel{3}578}
\l{}
\l{termografo. (\sou{fiziol.}). Instrumento qua konsistas ye aero-}
\l{\cel{3}termometro kun enrejistro-levero, ed uzata por mezurar}
\l{\cel{3}la varii di temperaturo, e la duro di ta varii, en irga}
\l{\cel{3}punto di la homo-korpo. - DEIR.}
\l{}
\l{\sou{termokemio}. (\sou{kemio)}.\cel{2}Fako di la kemio-cienco, qua traktas}
\l{\cel{3}la determino di la quanto de kaloro qua pleas rolo en}
\l{\cel{3}la kombini. -\cel{2}DEFIRS.}
\l{}
\l{\sou{termologio}. Fako di fiziko, qua traktas kaloro. - DEFIS.}
\l{}
\l{termometro. Aparato por la mezuro di la temperaturi, per la}
\l{\cel{3}dilato o la kontrakto di merkuro, di alkoholo, ec. en}
\l{\cel{3}tubo gradoza. - DEFIRS.}
\l{}
\l{\sou{termometrio}. (\sou{fiziko}).La mezurado di kaloro. - DEFIRS.}
\l{}
\l{termoskopo. (fiziko). Aparato por montrar la existo e la}
\l{\cel{3}signo di difero pri la temperaturi reciproka di du me-}
\l{\cel{3}dii, ma qua donas unike valori relativa pri la grandeso-}
\l{\cel{3}grado di ta diferi. - DEFIS.}
\l{}
\l{\sou{termostato}. (\sou{teknol.)}\cel{2}Ensemblo teknikala qua posibligas}
\l{\cel{3}la obteno, en klozajo (+kluzajo) di temperaturo konstanta.}
\l{}
\l{}
\l{\sou{termotaktismo}. (\sou{biol.}) Su-diplaso reaktala da protoplasmo,}
\l{\cel{3}influe da kaloro. - DEFIS.}
\l{}
\l{\sou{termoterapio}. (\sou{medic.)}\cel{3}Kuraco di la morbi per kaloro. -}
\l{\cel{3}DEFIS.}
\l{}
\l{\sou{termotropismo}.(\cel{1}\sou{biol.}) Influo, da kaloro, di la protoplas-}
\l{\cel{3}mo che la elementi tisua, e precipue di la direciono o}
\l{\cel{3}la sinso di la kresko che la tisui qui kreskas. - DEFIS.}
\l{}
\l{\sou{ternara}. (\sou{mat}.) (\sou{pri sistemo di nombrifado}). En qua tri}
\l{\cel{3}unaji ek la klaso unesma formacas unajo ek la klaso du-}
\l{\cel{3}esma. - DEFIS.}
\l{}
\l{\sou{terorar}. (\sou{netrans., pro, pri})\cel{2}Fremisar, parturbesar pro-}
\l{\cel{3}funde en sua organismo, pro pavoro. - DEFIS.}
\l{}
\l{\sou{terpentino}.- (\sou{kemio)}.Rezino liquida quan donas la arbori ek la}
\l{\cel{3}familio \textquotedbl{}terebintacei\textquotedbl{}, e di qua un elemento esas C10 H16}
\l{\cel{3}- DE.}
\l{}
\l{\sou{terpleno}. (\sou{militarto}).\cel{3}Platformo ek tero adjuntita, suste-}
\l{\cel{3}nata da parapeti masonura. - F.}
\l{}
\l{\sou{testo}. Persono qua deklaras (maxim-multa-kaze : koram judi-}
\l{\cel{3}ciisti) certifikar vidir od audir (ulo, ulu). - IS.}
\l{}
\l{testaceo. (zool.)\cel{2}Molusko di qua la korpo kovresas da kon-}
\l{\cel{3}ko. - DEFIS}
}
